\chapter*{Lời Mở Đầu}
\addcontentsline{toc}{chapter}{Lời Mở Đầu}
\markboth{Lời Mở Đầu}{Lời Mở Đầu}

\begin{truyen}{Lời Mở Đầu}{Opening Words}
\chuhoa{N}{hững câu nói ngắn của thầy giáo trong lớp học thường không được ghi vào sổ đầu bài, nhưng lại ở lại rất lâu trong ký ức học trò.}
\textit{(Short sentences a teacher says in class are rarely written in the lesson log, yet they stay in students' memory for a long time.)}

Có thể là một câu mở đầu tiết học, một lời kết khi hết giờ, hay một câu thầy viết lên bảng khi cả lớp đang im lặng suy nghĩ.
\textit{(It may be an opening line for the lesson, a closing thought when the bell rings, or a sentence the teacher writes on the board while the class is quietly thinking.)}

Những câu ấy không thay thế bài giảng, nhưng chúng làm cho việc học có hồn, và đôi khi chính chúng mới là thứ học sinh nhớ mãi.
\textit{(Those sentences do not replace the lesson, but they give learning a soul — and sometimes they are exactly what students remember forever.)}

Quyển này gom 50 câu như thế, chia làm năm nhóm: về việc học, về nỗ lực, về thái độ sống, về sai lầm, và về tương lai. Mỗi câu đi kèm một câu chuyện nhỏ, để thầy cô có thể dùng khi mở bài, kết tiết, viết bảng, hoặc tặng học sinh cuối năm.
\textit{(This volume gathers 50 such sentences in five groups: on learning, on effort, on attitude, on mistakes, and on the future. Each is paired with a short story, so teachers can use them to open a lesson, close a class, write on the board, or give to students at year's end.)}

Mong rằng mỗi trang sách sẽ là một món quà nhỏ — cho thầy cô thêm câu để nói, cho học trò thêm điều để giữ.
\textit{(May each page be a small gift: for teachers, one more sentence to say; for students, one more thing to keep.)}
\end{truyen}

\begin{baihoc}
Một câu nói đúng lúc có thể thay đổi cách một học sinh nghĩ về việc học và về chính mình.
\textit{(One sentence at the right moment can change how a student thinks about learning and about themselves.)}
\end{baihoc}

\begin{ghinhoanh}
\textbf{Short English to remember:}

A teacher's words in class can outlast the lesson.

\end{ghinhoanh}

\begin{tuhocnhanh}
1. Vì sao những câu nói ngắn của thầy lại đáng nhớ? \textit{Why do short teacher sayings stay memorable?}

2. Em nhớ nhất câu nào thầy cô từng nói trong lớp? \textit{Which sentence from your teacher do you remember most?}

3. Nếu em là thầy cô, em muốn học trò giữ lại câu gì? \textit{If you were the teacher, what would you want students to keep?}
\end{tuhocnhanh}

\ngancach
